\section{Đo trên ảnh thật}
\label{sec:ch10-anh-that}

\index{thiên kiến quy nạp!vì sao chương này cần dữ liệu thật}%
Chương~\ref{ch:tien-huan-luyen} không đo được luận điểm của nó, và mục 9.5 của
nó nói vì sao: corpus sinh từ văn phạm không mang cấu trúc mà luận điểm ấy nói
tới, và một bảng đo sai đại lượng tệ hơn không có bảng.

Chương này gặp cùng câu hỏi và trả lời khác, vì ở đây có đường ra. Luận điểm là
về ảnh, và ảnh thật thì tải về được.

Nên sách này tải CIFAR-10~\autocite{krizhevsky2009cifar}: 60,000 ảnh màu
32x32, 10 lớp, 50,000 ảnh huấn luyện và 10,000 ảnh kiểm tra, mỗi lớp chia đều.

Cái giá phải nói ra. Cho tới chương trước, mọi thí nghiệm trong repo companion
sinh dữ liệu của chúng từ một seed, nên \code{verify.py} chạy được trên máy
người khác mà không cần mạng. Từ chương này thì không. Lần chạy đầu phải tải
170,498,071 byte, và trên máy tôi viết chương này việc ấy mất ba mươi hai phút.
Đó không phải chi phí của bộ dữ liệu mà là của đường mạng ở đây, nhưng nó lớn
tới mức không nhét vào ngân sách nào được, nên nó nằm ngoài phần bấm giờ: tải
một lần rồi mọi lần sau đọc cache. Một người đọc không có mạng sẽ thấy mục này
đỏ cho tới khi tải xong một lần.

Trang tải ghi con số ấy là \enquote{163 MB}, tức nó đang gọi MiB là MB; tôi in
số byte vì đó là số \code{cifar.py} đối chiếu trước khi giải nén.

Tôi đổi tính chất ấy lấy một chương có bằng chứng, và nói ra ở đây để bạn biết
mình đang trả gì.

\subsection*{Ba mô hình, và số tham số nói trước}

\index{thiên kiến quy nạp!ba mô hình đem so}%
Phép so là một mạng \tn{tích chập}{convolution} nhỏ đặt cạnh những mạng
perceptron nhiều tầng - multi-layer perceptron, viết tắt là MLP - nhìn cùng số
pixel ấy mà không có \tn{tính cục bộ}{locality} và không có
\tn{chia sẻ trọng số}{weight sharing}.

\begin{measured}{ba mô hình của mục này}
\begin{minted}{text}
model         parameters  vs CNN  locality  sharing
CNN           66,570      1.0000  yes       yes
MLP matched   64,753      0.9727  no        no
MLP wide      1,578,506   23.7120 no        no
\end{minted}

Đo bằng \code{experiments/ch10\_cifar.py} tại tag \code{ch10}.
\end{measured}

Quyết định của sách là một \tn{phép cắt bỏ}{ablation} phải có đối chứng khớp
kích thước hoặc nói rõ vì sao không có. Ở đây có hai đối chứng, và cái thứ hai mới là cái đóng
lập luận.

MLP khớp tham số có \textbf{21 unit ẩn}. Đó là thứ 66,570 tham số mua được khi
không tham số nào dùng chung, và nó là hệ quả trực tiếp của con số 1024 ở
mục~\ref{sec:ch10-hai-rang-buoc}: mỗi unit ẩn phải trả riêng cho từng pixel
trong 3072 pixel. Chỉ so với nó thôi thì phản bác hiển nhiên là \enquote{21
unit thì làm được gì}, và phản bác ấy đúng. Nên bảng có thêm một MLP với 23.71
lần số tham số.

\subsection*{Bao nhiêu ảnh thì đủ}

\index{thiên kiến quy nạp!bảng sample efficiency trên CIFAR-10}%
Cùng một recipe cho mọi hàng: Adam, learning rate \(10^{-3}\), batch 128. Cái
thay đổi giữa các hàng là lượng dữ liệu. Ba seed, chấm trên toàn bộ 10,000 ảnh
kiểm tra.

\begin{measured}{độ chính xác theo lượng dữ liệu huấn luyện}
\begin{minted}{text}
n_train  epochs  model         test acc     spread   train acc
1000     25      CNN           0.4365       0.0326   0.8893
1000     25      MLP matched   0.3163       0.0176   0.9190
1000     25      MLP wide      0.3469       0.0236   1.0000

4000     14      CNN           0.5499       0.0102   0.7637
4000     14      MLP matched   0.3733       0.0042   0.6432
4000     14      MLP wide      0.3915       0.0133   0.9453

16000    7       CNN           0.6399       0.0197   0.7202
16000    7       MLP matched   0.4251       0.0152   0.5100
16000    7       MLP wide      0.4534       0.0120   0.6835

50000    4       CNN           0.7036       0.0019   0.7361
50000    4       MLP matched   0.4478       0.0015   0.4797
50000    4       MLP wide      0.5025       0.0057   0.5743
\end{minted}

Số epoch giảm theo lượng dữ liệu vì 1000 ảnh với batch 128 chỉ là 8 bước
gradient một epoch; giữ epoch cố định là so 8 bước với 390 bước.

Đo bằng \code{experiments/ch10\_cifar.py} tại tag \code{ch10}.
\end{measured}

Cột \code{train acc} có mặt để phân biệt hai chuyện khác nhau: một mô hình
không khớp nổi dữ liệu, và một mô hình khớp được mà không tổng quát hóa. Ở hàng
1000 ảnh, cả ba mô hình đều khớp gần hết tập huấn luyện - MLP wide khớp đúng
1.0000 - nên không mô hình nào ở đó đang thiếu bước. Đó là phép kiểm rẻ cho một
cái bẫy tôi đã sập một lần trong chính phiên này; mục~\ref{sec:ch10-con-lai-gi}
kể.

\begin{measured}{khoảng cách, đặt cạnh spread của seed}
\begin{minted}{text}
n_train  CNN-matched  CNN-wide  CNN spread  widest MLP spread
1000     0.1202       0.0897    0.0326      0.0236
4000     0.1767       0.1584    0.0102      0.0133
16000    0.2148       0.1865    0.0197      0.0152
50000    0.2558       0.2011    0.0019      0.0057
\end{minted}

Đo bằng \code{experiments/ch10\_cifar.py} tại tag \code{ch10}.
\end{measured}

Đọc cột khoảng cách cạnh cột spread, vì chương~\ref{ch:chi-phi-va-song-song} là
chỗ sách này học rằng một khoảng cách hẹp hơn spread thì không phải kết quả.
Chỗ chặt nhất trong bảng là ô hẹp nhất của cả hai cột khoảng cách, 0.0897 ở
hàng 1000 ảnh, đặt cạnh spread lớn nhất của chính hàng ấy, 0.0326: rộng gấp
2.75 lần. Không phải một tỷ lệ áp đảo, và nó là hàng ít dữ liệu nhất nên cũng
là hàng nhiễu nhất; các hàng còn lại rộng hơn nhiều. Không hàng nào lật chiều
được, nhưng hàng 1000 ảnh là hàng tôi sẽ không xây một lập luận mảnh lên trên.

\subsection*{Chỗ tôi suýt viết ngược}

\index{thiên kiến quy nạp!khoảng cách rộng ra chứ không hẹp lại}%
Câu tôi định viết trước khi chạy là câu quen thuộc:
\tn{thiên kiến quy nạp}{inductive bias} đáng giá nhất khi ít dữ liệu, nên
khoảng cách phải \emph{hẹp lại} khi dữ liệu nhiều lên.

Bảng nói ngược. Khoảng cách rộng ra, 0.1202 lên 0.2558.

Chỗ nhầm không nằm ở số mà ở đại lượng. \enquote{Khoảng cách ở một lượng dữ
liệu cố định} và \enquote{số ảnh cần để đạt cùng độ chính xác} là hai câu hỏi
khác nhau, và chỉ câu thứ hai nói về việc dùng dữ liệu hiệu quả tới đâu. Hỏi
câu thứ hai thì được cái này:

\begin{measured}{số ảnh mỗi MLP cần để đạt cái CNN đạt}
\begin{minted}{text}
CNN trained on  reaches  MLP matched needs  MLP wide needs
1000            0.4365       28,403 (28.4x)     10,958 (11.0x)
4000            0.5499          over 50,000        over 50,000
16000           0.6399          over 50,000        over 50,000
50000           0.7036          over 50,000        over 50,000
\end{minted}

Số ảnh cần là nội suy log-tuyến tính giữa hai hàng kề nhau.

Đo bằng \code{experiments/ch10\_cifar.py} tại tag \code{ch10}.
\end{measured}

\textbf{28.4 lần.} MLP khớp tham số cần 28,403 ảnh để đạt cái CNN đạt được
bằng 1000. Đó là tỷ giá mà mục~\ref{sec:ch10-attention-lam-duoc} nói tới, đo ở
quy mô này: hai ràng buộc, nướng sẵn vào tầng, đáng gần một bậc rưỡi độ lớn về
dữ liệu.

Ba hàng còn lại ghi \enquote{over 50,000} chứ không ngoại suy. Từ 4000 ảnh trở
lên, không lượng dữ liệu nào \emph{trong bộ dữ liệu này} đưa MLP tới chỗ CNN
đạt được, và ngoại suy ra ngoài dải đã đo là đúng chuyện
chương~\ref{ch:tien-huan-luyen} đã dành một bài tập để nói.

Cột cuối là cột trả lời phản bác. Cho MLP 23.71 lần tham số thì nó vẫn cần
10,958 ảnh, tức 11.0 lần. Bội số nhỏ đi rõ rệt và \emph{không} biến mất: đổ
thêm tham số mua lại được hơn một nửa khoảng cách và không mua hết.

Còn vì sao khoảng cách tuyệt đối rộng ra thì cũng đọc được từ chính bảng trước.
MLP khớp tham số bão hòa sớm vì nó chỉ có 21 unit ẩn: nó bị chặn bởi sức chứa
chứ không bởi dữ liệu. Đó đúng là lý do bảng phải có hàng MLP wide.

\subsection*{Mô hình đã huấn luyện có bất biến với phép dịch không}

\index{thiên kiến quy nạp!đo bất biến trên mô hình đã huấn luyện}%
Mục~\ref{sec:ch10-dang-bien} đo \tn{đẳng biến}{equivariance} tại khởi tạo, trên
một tầng. Câu hỏi
còn lại là câu người ta thật sự muốn hỏi: một mô hình đã học xong có nhận ra
cùng một ảnh khi ảnh dịch đi không.

\begin{measured}{tỷ lệ ảnh kiểm tra đổi nhãn dự đoán khi dịch}
\begin{minted}{text}
shift  CNN     MLP matched  MLP wide
1      0.1145  0.2409       0.2158
2      0.1817  0.3790       0.3478
4      0.3078  0.5446       0.5106
8      0.5981  0.7343       0.7013
\end{minted}

Phép dịch ở đây là dịch vòng, nên dịch 8 pixel trên ảnh 32 pixel chuyển hẳn
một phần tư ảnh sang phía bên kia.

Đo bằng \code{experiments/ch10\_cifar.py} tại tag \code{ch10}.
\end{measured}

Phép so đáng đọc là theo cột, cùng một phép dịch. Ở dịch một pixel, CNN đổi ý
trên 11.45\% số ảnh còn MLP khớp tham số đổi trên 24.09\%: chưa bằng một nửa.
Thiên kiến quy nạp mua được cái đó, và nó mua được thật.

Nhưng nhìn hàng cuối. Ở phép dịch 8, CNN vẫn đổi ý trên 59.81\% số ảnh, và
\emph{không} có mâu thuẫn nào với con số không tuyệt đối ở
mục~\ref{sec:ch10-dang-bien}. Có ba lý do rời nhau, và đáng tách ra vì trộn
chúng lại là cách người ta đi tới câu \enquote{mạng tích chập bất biến với phép
dịch}.

Thứ nhất, mạng này pad bằng không chứ không pad vòng. Phép đo cho số không
tuyệt đối ở mục trước dùng padding vòng; với padding không thì
mục~\ref{sec:ch10-dang-bien} đã đo được rằng đẳng biến hỏng trên khung viền. Mà
phép dịch vòng ở đây chuyển nội dung qua đúng cái viền ấy.

Thứ hai, kể cả nếu đặc trưng đẳng biến tuyệt đối thì đầu ra vẫn không bất biến.
Đẳng biến nói \tn{bản đồ đặc trưng}{feature map} dịch theo. Bộ phân loại đứng
sau là một tầng Linear đọc bản đồ ấy đã phẳng hóa, nên bản đồ dịch đi một ô là
một vector khác hẳn, và nó có toàn quyền trả lời khác. Đẳng biến của đặc trưng
không kéo theo \tn{bất biến}{invariance} của mạng, và không có gì trong kiến
trúc này bắt nó phải kéo theo.

Thứ ba, dịch vòng 8 pixel thật sự đổi ảnh. Một phần tư ảnh sang phía bên kia là
một ảnh khác, không phải cùng ảnh nhìn lệch đi.

Tôi không tách được ba phần ấy bằng bảng này, và bài tập 8 tier một giao lại
việc thiết kế phép đo tách được. Cái bảng này nói chắc chắn là phép so theo
cột, và nó đủ cho luận điểm.
